\section{Legal Information Retrieval}

\subsection{Characteristics and Challenges}

The process of legal research is an integral component of a legal professional's daily workflow and the development of tool support for this process, known as \acl{LIR} (\acs{LIR}), has been a cornerstone of \ac{AI} \& Law research \cite{van_opijnen_concept_2017}\citep[p. 21]{landthaler_improving_2019}. \citet{van_opijnen_concept_2017} provide a foundational theoretical framework for \ac{LIR}, which is defined as a specific subtype of \ac{IR} that has been justified by a variety of characteristics particular to legal text data. These include among other things specialised terminology, extensive synonymy, dense cross-referencing and a diversity of document types\footnote{A more detailed definition of the legal domain can be found in Section \ref{sec:domain-identification}} \cite{van_opijnen_concept_2017}. They designed a six-dimensional model of relevance in \ac{LIR} that extends beyond simple topical matching which highlights the unique structure and interconnectedness of legal information objects \cite{van_opijnen_concept_2017}.

Most research in \ac{LIR} addresses challenges that arise from the specific domain characteristics \cite{hafner_1978, bing_designing_1987, Dick1991_Representation, Ashley_2017}. A major trend has been to leverage \ac{NLP} technologies to access deeper semantic layers of legal documents. For instance, \citet{Grabmair_2015} combined an ontology of intellectual property law with \ac{NLP} and machine learning to improve conceptual legal information retrieval of vaccine injury decisions. \citet{schweighofer2007legal} explored ontologies and relevance feedback methods for query expansion. \citet{waltl2018semantic} engineered a sophisticated \ac{NLP}-based tool suite for legal experts that includes both basic and advanced legal information retrieval functionality using the UIMA type system. Empirical studies emphasize the practical importance of good legal search or retrival. \citet{lastres2015rebooting}, surveying 190 law firm associates, found that up to 15 hours per week are spent conducting legal research, while user studies by \citet{peoples2005death} and \citet{wiggers2018exploration} investigated how legal professionals interact with different search modalities and relevance criteria. 

The general challenges of the legal domain are fully reflected in the German legal domain. As \citet{landthaler_improving_2019} observes, German legal information retrieval constitutes a rarely investigated topic in research, despite the distinctive challenges it poses and more research on German \ac{AI} \& Law is required \cite{waltl2018semantic}. German tax law in particular is widely regarded as an very complex legal field. %(src: 4h ps://www.spiegel.de/wirtschaft/service/steuern-kommt-die-mehrheit-der-weltweiten-steuerliteratur-ausdeutschland-a-1111192.html, last accessed August 10, 2019 5h ps://www.deutscheranwaltspiegel.de/auf-die-positionierung-kommt-es-an/, last accessed August 10, 2019)

\subsection{Retrieval Approaches for German Legal Text}

The initial approaches to \ac{LIR} were largely based on traditional term-based vector space models such as \ac{TF-IDF} and Okapi BM25, which capture lexical overlap but struggle to account for the pervasive synonymy and semantic complexity of legal language \cite{wiggers2018exploration}. To supplement these purely term-based methods, early advanced approaches explored the extraction and exploitation of citations and cross-references within legal documents as a complementary retrieval signal \cite{Waltl2016_Differentiation,landthaler_improving_2019}.

With the introduction of distributional semantic embedding-models, multiple studies have explored their appllication to legal text representation. For example, \citet{sugathadasa2019legal} compared \ac{TF-IDF} with Paragraph Vectors (doc2vec) for legal document retrieval. \citet{adebayo2016textual} applied normalised word embedding vectors to \ac{IR} and textual entailment tasks on the COLIEE competition dataset. \citet{vo2017experimenting} used word embeddings to suggest synonyms during technology-assisted legal review and \citet{swyminski2019automatic} constructed a legal dictionary for Polish law using word embeddings. Consistently, these studies demonstrate the superiority of embedding-based approaches over pure keyword matching in \ac{LIR}-tasks. \citet{landthaler_improving_2019} confirms this finding for German legal search specifically, showing that static word embeddings effectively capture synonymy relations for query expansion in German tax law. He also proposed the WE-DF vector space model as an implicit query expansion mechanism that eliminates the need for manually maintained thesaurian approach, similarly explored by \citet{pickel2016prototypical} on the German Civil Code. However, those contributions operate entirely within the pre-Transformer paradigm of static, context-independent word embeddings, leaving modern contextualised document representations unexplored.

\subsection{Evaluating Retrieval in Legal \ac{RAG} Systems}

Since finding the most relevant documents from a large corpus given a user query is indipendet of whether the retrieved documents are presented directly to a user or passed as context to a generative model, the requirements for the retriever component of an \ac{RAG} application are identical to those in the classic \ac{LIR}. Therefore, the domain-specific challenges identified in the \ac{LIR} literature apply in full to the retrieval stage of a \ac{RAG} system operating on legal text \cite{van_opijnen_concept_2017}\citep[p. 24]{landthaler_improving_2019}.

Recently, efforts have been made to isolate and address this specific retrieval component of \ac{RAG} systems in legal settings through dedicated evaluations. \citet{pipitone2024legalbench} introduced LegalBench-\ac{RAG}, the first benchmark explicitly designed to evaluate retrieval quality in legal \ac{RAG} independently of the downstream generator. Building on this benchmark, \citet{reuter_towards_2025} identified a critical failure mode called \ac{DRM}, where retrievers select chunks from entirely incorrect source documents due to high lexical and structural similarity within legal corpora, observing \ac{DRM} rates exceeding 95\% in standard \ac{RAG} configurations. Their proposed mitigation, Summary-Augmented Chunking, enriches each text chunk with a document-level synthetic summary to preserve global context. While their work demonstrates the vulnerability of retrieval in legal domains, it is restricted to English-language data and focuses strictly on the pre-retrieval pipeline, rather than the representational deficiencies of the embedding models themselves\footnote{While a vast body of literature focuses on architectural \ac{RAG} optimisations, such as query expansion, hybrid search methodologies, or neural re-ranking, these techniques often act as compensatory layers that mask the representational deficiencies of the base embeddings \cite{gao_2024_rag}. Therefore, investigating the root cause within the embedding space itself remains a critical, yet under-explored, necessity.} \cite{reuter_towards_2025}. 

Contrary, for the German legal domain,  evaluation has traditionally relied on classical \ac{IR} benchmarks rather than \ac{RAG}-specific frameworks. \citet{wrzalik_gerdalir_2021} introduced \ac{GerDaLIR}, the first large-scale German legal information retrieval dataset, comprising over 131,000 collection documents segmented into more than 3 million passages with approximately 123,000 queries derived from case law references. Their baseline experiments compare term-based methods, word-centroid similarity using word embeddings and neural reranking with German \ac{BERT}  and ELECTRA models fine-tuned on the dataset. However, it covers general German legal language, predominantly civil and administrative law, rather than the highly specialised tax domain.  Building on \ac{GerDaLIR}, \citet{bönisch2025finding} proposed an SVR-ensemble approach over Transformer-based embedding spaces and observed that the embedding space topology is heavily influenced by passage length, with texts clustering together based on length regardless of semantic content. However, their focus remains on the downstream retrieval methodology, the ensemble classification layer, rather than on evaluating modern embedding models or diagnosing and improving the intrinsic quality of the embedding space itself.